\article{$\theta$ 级数与不变量理论}{$\theta$-series and invariant theory}{R. Howe}{1}{275--285}
\origsection{0}{引言（Introduction）}
在自守形式理论的发展中，$\theta$ 级数（theta-series）一直是构造自守形式（automorphic form）的主要方法之一。在 [Wi 1] 中，Weil 指出，一个在量子力学中也有重要作用的酉表示（unitary representation）构成了 $\theta$ 级数理论的基础，尤其是 Siegel [Si 1]、[Si 2] 所发展的理论。本文概述以表示论语言表述并推广 Weil 的 $\theta$ 级数理论的工作。到目前为止，这一表述提供了一个系统框架，将若干作者近期的工作与较早的一些重要计算统一起来（例子见 [G]）。这种方法的一个关键特点，是它与 [Wy] 所阐述的经典不变量理论（classical invariant theory）有密切联系。事实上，$\theta$ 级数理论在某种意义上似乎是经典不变量理论的超越版本。
\paragraph*{约定（Convention）}
全文中 $F$ 表示一个域，特征始终不等于 $2$。对 $F$ 的进一步假设随讨论内容而变，并在相应段落开头说明。
\origsection{1}{Heisenberg 群与 Stone--von Neumann 定理：局部情形}
设 $F$ 为局部域（local field），$W$ 是 $F$ 上的辛向量空间（symplectic vector space），其形式为 $\langle\ ,\ \rangle$。定义相应的 Heisenberg 群（Heisenberg group）$H(W)$：作为集合，$H(W)=W\oplus F$，群运算为
\[
(w,s)(w',s')=(w+w',s+s'+\tfrac12\langle w,w'\rangle).
\]
令 $\chi$ 为 $F$ 的非平凡特征（nontrivial character）。
\begin{originaltheorem}{Stone--von Neumann 定理（Stone--von Neumann theorem）}
中心特征（central character）为 $\chi$ 的 $H$ 的不可约酉表示，只有一个等价类。
\end{originaltheorem}
用 $\rho_\chi$ 表示这一表示；其在光滑向量（smooth vector）上的相应表示 [C 2]、[Wa] 记作 $\rho_\chi^\infty$。

这一表示的等价类虽然唯一，却有许多具体实现；正如 Cartier [C 1] 所指出的，这些实现可以组织成光滑族。这种现象体现了振子表示理论的一个核心特点。
\origsection{2}{振子表示：局部情形（The oscillator representation, local case）}
$F$ 仍为局部域。令 $\operatorname{Sp}$ 为 $\langle\ ,\ \rangle$ 的等距群（isometry group）。对 $g\in\operatorname{Sp}$ 和 $(w,s)\in H(W)$，作用 $g((w,s))=(g(w),s)$ 将 $\operatorname{Sp}$ 嵌入 $H(W)$ 的自同构群。

根据 [M]，$\operatorname{Sp}$ 有唯一的非平凡二重覆盖（two-fold cover）$\widetilde{\operatorname{Sp}}$（$F=\C$ 的情形除外，此时约定 $\widetilde{\operatorname{Sp}}=\operatorname{Sp}$）。以 $\widetilde g\mapsto g$ 表示覆盖映射。
\begin{originaltheorem}{定理（Shale--Weil）}
在 $\rho_\chi$ 的空间上存在 $\widetilde{\operatorname{Sp}}$ 的表示 $\omega_\chi$（在酉等价意义下唯一），使对 $\widetilde g\in\widetilde{\operatorname{Sp}}$、$h\in H$ 有
\[
\omega_\chi(\widetilde g)\rho_\chi(h)\omega_\chi(\widetilde g^{-1})=\rho_\chi(g(h)).\tag{2.1}
\]
\end{originaltheorem}
$\omega_\chi$ 有相应的光滑表示 $\omega_\chi^\infty$。$\rho_\chi$ 和 $\omega_\chi$ 的光滑向量相同（在任意共同实现中），所以 $\rho_\chi^\infty$ 和 $\omega_\chi^\infty$ 作用在同一空间上。

称 $\omega_\chi$，或 $\operatorname{Sp}$ 的相应乘子表示（multiplier representation），为振子表示（oscillator representation）。
\origsection{3}{阿代尔化（Adelization）}
现在令 $F$ 为整体域（global field），也称 $A$ 域。以 $v$ 表示 $F$ 的素位，$F_v$ 为 $F$ 在 $v$ 处的完备化（completion），$\A$ 为 $F$ 的阿代尔环（adele ring）。若 $W=W_F$ 是 $F$ 上的辛向量空间，可构造 $H(W_F)$ 及其阿代尔化 $H(W_\A)$。照例，$H(W_F)$ 作为离散子群嵌入 $H(W_\A)$，且商空间紧。下文省略 $W$，简记 $H(W_F)=H_F$、$H(W_{F_v})=H_v$ 等。

设 $\{e_i,f_j\}$ 为 $W$ 的辛基（symplectic basis）。令 $\Gamma_v$ 为 $H_F$ 中元素 $\{(e_i,0),(f_j,0)\}$ 所生成子群在 $H_v$ 中的闭包。那么 $H_\A$ 是 $H_v$ 关于 $\Gamma_v$ 的限制直积（restricted direct product）。\footnote{\textbf{译者注（本次校订）：}对于一般整体域，仅由所列有限个基向量作群生成，再取闭包，未必得到紧开子群。例如 $F_v/\Q_p$ 的次数大于一时，这种闭包只有 $\Z_p$ 系数，不能产生完整的 $\mathcal O_v$ 格。应从整数格出发，在几乎所有素位取相应的紧开 Heisenberg 子群；当 $v\nmid2$ 时，可用辛格 $L_v$ 构造 $L_v\oplus\mathcal O_v$。}

令 $\chi$ 为 $\A$ 的一个非平凡特征，且 $F\subseteq\ker\chi$；令 $\{\chi_v\}$ 为 $\chi$ 的局部因子。在几乎所有素位，都有唯一的被 $\rho_{\chi_v}(\Gamma_v)$ 固定的向量 $\phi_v$。\footnote{\textbf{译者注（本次校订）：}这里的“唯一”应理解为不变向量空间是一维的，即非零向量在相差标量的意义下唯一；构造限制张量积时须从中选定一个非零向量。下文对 $J_v$ 不变向量的同类说法也按此理解。}因此可以取 $\rho_{\chi_v}$ 的限制张量积（restricted tensor product），构成 $H(W_\A)$ 的表示 $\rho_\chi$。由于每个表示 $\rho_{\chi_v}^\infty$ 都是容许的（admissible）[C 2]，每个 $\phi_v$ 都是光滑向量，$\rho_\chi$ 的光滑子表示 $\rho_\chi^\infty$ 是容许的，并且是 $\rho_{\chi_v}^\infty$ 的限制张量积 [F]。

令 $J_v\subseteq\operatorname{Sp}_v$ 为 $\Gamma_v$ 的正规化子（normalizer）。那么 $\operatorname{Sp}_\A$ 是 $\operatorname{Sp}_v$ 关于 $J_v$ 的限制直积。将各局部作用组合起来，可得到 $\operatorname{Sp}_\A$ 在 $H_\A$ 上的作用。如同局部情形，存在二重覆盖 $\widetilde{\operatorname{Sp}}_\A$ 及其表示 $\omega_\chi$，使 (2.1) 的整体对应式成立。存在相应的光滑表示 $\omega_\chi^\infty$。尽管 $\widetilde{\operatorname{Sp}}_\A$ 不是 $\widetilde{\operatorname{Sp}}_v$ 的限制直积，仍有交换图（commutative diagram）
\[
\begin{tikzcd}
\prod_v\widetilde{\operatorname{Sp}}_v\arrow[r]\arrow[d]&\widetilde{\operatorname{Sp}}_\A\arrow[d]\\
\prod_v\operatorname{Sp}_v\arrow[r]&\operatorname{Sp}_\A
\end{tikzcd}\tag{3.1}
\]
(3.1) 中的积均为关于 $J_v$ 的限制直积。（对几乎所有 $v$，事实上对 $v\nmid2$，存在 $J_v$ 到 $\widetilde{\operatorname{Sp}}_v$ 的唯一提升。方便时利用这一提升，将 $J_v$ 看作 $\widetilde{\operatorname{Sp}}_v$ 的子群。）对几乎所有 $v$，向量 $\phi_v$ 是被 $\omega_{\chi_v}^\infty(J_v)$ 固定的唯一向量，且 $\omega_\chi^\infty$ 到 $\prod_v\widetilde{\operatorname{Sp}}_v$ 的拉回是 $\omega_{\chi_v}^\infty$ 关于 $\phi_v$ 的限制张量积。
\begin{originaltheorem}{定理（Weil）}
存在唯一同态 $\beta:\operatorname{Sp}_F\to\widetilde{\operatorname{Sp}}_\A$，使三角图
\[
\begin{tikzcd}
&\widetilde{\operatorname{Sp}}_\A\arrow[d]\\
\operatorname{Sp}_F\arrow[ur,"\beta"]\arrow[r]&\operatorname{Sp}_\A
\end{tikzcd}
\]
交换，其中水平箭头与竖直箭头分别为典范单射与投影。
\end{originaltheorem}
因此对 $g\in\operatorname{Sp}_F$，可以无歧义地谈论 $\omega_\chi(g)$。
\origsection{4}{$\Theta$ 分布（The theta-distribution）}
本节继续采用第3节对 $F$ 的约定。
\begin{originaltheorem}{定理 4.1（Theorem 4.1）}
在 $\rho_\chi^\infty$（的空间 $Y$）上存在一个在 $H_F$ 下不变的线性泛函（linear functional）$\Theta$。即对 $h\in H_F,y\in Y$，有 $\Theta(\rho_\chi^\infty(h)(y))=\Theta(y)$。泛函 $\Theta$ 在相差常数倍的意义下唯一。
\end{originaltheorem}
这一定理本质上是 Poisson 求和公式（Poisson summation formula）的一个强形式。[C 1] 中有同类结果。$\Theta$ 的存在隐含于 Weil [Wi 1] 中，并在经典理论里得到运用。第3节中提升 $\beta$ 的构造，也主要依赖于 $\Theta$ 的存在。

由于 $\operatorname{Sp}_F$ 正规化 $H_F$，由 (2.1) 和 $\Theta$ 的唯一性可见，$\omega_\chi(\operatorname{Sp}_F)$ 在相差常数倍的意义下保持 $\Theta$。因为 $\operatorname{Sp}_F$ 等于自身的交换子群，$\Theta$ 实际上在 $\operatorname{Sp}_F$ 下不变。因此，若 $u$ 为 $\rho_\chi^\infty$ 空间中的向量，则 $\widetilde{\operatorname{Sp}}_\A$ 上的函数
\[
\theta(u)(g)=\Theta(\omega_\chi(g)u),\qquad g\in\widetilde{\operatorname{Sp}}_\A,
\]
可下降为陪集空间 $X=\operatorname{Sp}_F\backslash\widetilde{\operatorname{Sp}}_\A$上的函数。还可验证，$\theta(u)(g)$ 在 $X$ 的无穷远处具有“适度增长”（moderate growth），因此 $\theta(u)$ 实际上是 $\widetilde{\operatorname{Sp}}_\A$ 上的自守形式。称 $\theta(u)$ 为与 $u$ 相联系的 theta 级数。

若 $\dim W=2$，且适当选取 $u$，则 $\theta(u)$ 就是 Jacobi 的经典 $\theta$ 级数。适当放宽 $u$ 的选择范围，则得到“带特征的 $\theta$ 常数”（theta-constants with characteristic）[I]。因此，$\theta$ 分布产生推广了熟知的 $\theta$ 级数的自守形式。这些 $\theta(u)$ 仍是一类十分特殊的自守形式。进一步，若 $G\subseteq\operatorname{Sp}_F$ 是约化代数子群（reductive algebraic subgroup），则可将 $\theta(u)$ 限制到 $\widetilde G_\A$（它是 $G_\A\subseteq\operatorname{Sp}_\A$ 在 $\widetilde{\operatorname{Sp}}_\A$ 中的逆像），从而得到 $\widetilde G_\A$ 上的自守形式。若 $G$ 是相对较小的子群，便能得到更广泛的 $\widetilde G_\A$ 上的自守形式。对于一般的 $G$，很难十分精确地说明 $\theta(u)|_{\widetilde G_\A}$ 的性质；但若 $G$ 属于下面定义的约化对偶对，则预期可以较为详细地描述其结构。$\theta$ 级数与自守形式理论的关联，在这一情形下表现得尤为清楚。
\origsection{5}{约化对偶对（Reductive dual pairs）}
这里 $F$ 任意（但特征仍不等于 $2$）。设 $(G,G')$ 为 $\operatorname{Sp}$ 的一对子群。若满足
\begin{enumerate}
\item[(i)] $G,G'$ 在 $W$ 上绝对约化地（absolutely reductively）作用；
\item[(ii)] $G$ 是 $G'$ 在 $\operatorname{Sp}$ 中的中心化子（centralizer），反之亦然，
\end{enumerate}
则称 $(G,G')$ 构成约化对偶对。

约化对偶对的分类如下。若 $(G,G')$ 是 $\operatorname{Sp}$ 中的约化对偶对，而 $W=W_1\oplus W_2$ 是正交直和分解，$W_1,W_2$ 在 $G\cdot G'$ 下不变，则称 $(G,G')$ 可约（reducible）。$(G,G')$ 到 $W_i$ 的限制 $(G_i,G_i')$ 定义 $\operatorname{Sp}(W_i)$ 中的约化对偶对。称 $(G,G')$ 为 $(G_i,G_i')$ 的直和。若 $(G,G')$ 不可约分，则称为不可约（irreducible）。任意对偶对都能分解为不可约对偶对的直和，且这一分解在自然的等价意义下唯一。不可约对可如下描述。

\paragraph*{II 型（Type II）}
存在极大迷向子空间（maximal isotropic subspace）$X,Y\subseteq W$，使 $X\oplus Y=W$，并且 $X,Y$ 在 $G\cdot G'$ 下不变。这时存在：
\begin{enumerate}
\item[(a)] $F$ 上的除代数（division algebra）$D$（未必是中心的）；
\item[(b)] 右 $D$ 模 $X_1$ 及左 $D$ 模 $X_2$；使得
\item[(c)] $X\cong X_1\otimes_D X_2$，并在此认同下 $(G,G')$ 对应于 $(\GL_D(X_1),\GL_D(X_2))$。
\end{enumerate}
\paragraph*{I 型（Type I）}
$G\cdot G'$ 在 $W$ 上的联合作用不可约。这时存在
\begin{enumerate}
\item[(a)] 一个除代数 $D$；
\item[(b)] 其上的对合（involution）$\natural$；以及
\item[(c)] $D$ 模 $V_1,V_2$；
\item[(d)] 配备形式 $(\ ,\ )_1,(\ ,\ )_2$，一个是 $\natural$ Hermitian 形式，另一个是 $\natural$ 斜 Hermitian 形式（skew-Hermitian form）；使得
\item[(e)] $W\cong V_1\otimes V_2$，并在此认同下 $(G,G')$ 分别对应于 $(\ ,\ )_1,(\ ,\ )_2$ 的等距群；且
\item[(f)] $\langle\ ,\ \rangle=\operatorname{tr}_{D/F}((\ ,\ )_1\otimes(\ ,\ )_2^\natural)$。这里 $\operatorname{tr}_{D/F}$ 为约化迹（reduced trace）。代数 $D$ 未必在 $F$ 上是中心的。这里先在纯张量上按所给公式计算，再取迹；所得表达式与张量的表示方式无关，因而定义了 $F$ 双线性形式。
\end{enumerate}
\paragraph*{注（Remarks）}
（a）本质相同的分类可见于 Weil [Wi 2] 及别处，基本上可追溯至 Albert，也可由 [Sa] 的结果推出。

（b）局部域 $F$ 上带对合的除代数，或者是某个 $F$ 扩域上的二次扩张，或者是该扩域上的四元数代数（quaternion algebra）。\footnote{\textbf{译者注（本次校订）：}此处还应包括带恒等对合的交换域，例如 $D=F$；这正是正交群与辛群对偶对所用的情形。此外，当 $D$ 的中心 $E$ 严格大于 $F$ 时，上文的 $\operatorname{tr}_{D/F}$ 应理解为 $\operatorname{Tr}_{E/F}\circ\operatorname{Trd}_{D/E}$。}

（c）考虑数域（number field）$F$ 上的对 $(G,G')$。假设它是 I 型不可约对，$\natural$ 是 $D$ 上的正对合（positive involution），$V_2\cong D$ 配备 $\natural$ Hermitian 形式 $(x,y)_2=xy^\natural$。则 $V_1\cong W$ 带有 $\natural$ 斜 Hermitian 形式 $(\ ,\ )_1$。数据 $D,\natural,V_1,(\ ,\ )_2$ 是 Shimura [Sh] 用于指定其 PEL 型（PEL-type）的七项数据中的四项。Shimura 的其余数据出现在约化对偶对理论的进一步发展中。因此，Shimura 的 PEL 型由此自然地进入约化对偶对的理论。一个值得研究的问题是：Shimura 及其他人关于这些对象的结果，有多少可以从振子表示出发得到解释。
\origsection{6}{局部对偶性（Local duality）}
$F$ 现在又是局部域。设 $H\subseteq\widetilde{\operatorname{Sp}}$ 是闭子群，定义
\[
\mathcal R(H)=\left\{\begin{array}{l|l}
\sigma&\sigma\text{ 是 }H\text{ 的不可约光滑表示，且存在非平凡的}\\
&H\text{ 交织映射 }\alpha:\omega_\chi^\infty\longrightarrow\sigma
\end{array}\right\}.
\]
如前，若 $G\subseteq\operatorname{Sp}$，则 $\widetilde G$ 是 $G$ 在 $\widetilde{\operatorname{Sp}}$ 中的逆像。若 $(G,G')\subseteq\operatorname{Sp}$ 是约化对偶对，则 $\widetilde G,\widetilde G'$ 在 $\widetilde{\operatorname{Sp}}$ 中交换，所以 $\widetilde G\cdot\widetilde G'$ 的一个不可约容许表示拉回为 $\widetilde G\times\widetilde G'$ 的不可约容许表示。这样的表示形如 $\sigma\otimes\sigma'$，其中 $\sigma,\sigma'$ 分别是 $\widetilde G,\widetilde G'$ 的不可约容许表示。因此 $\mathcal R(\widetilde G\cdot\widetilde G')$ 定义了 $\widetilde G$ 的表示与 $\widetilde G'$ 的表示之间的对应（即给出这一对应的图像）。
\begin{originaltheorem}{猜想：局部对偶性（Conjecture: local duality）}
若 $(G,G')$ 是 $\operatorname{Sp}$ 中的约化对偶对，则 $\mathcal R(\widetilde G\cdot\widetilde G')$ 是 $\mathcal R(\widetilde G)$ 与 $\mathcal R(\widetilde G')$ 之间一个双射的图像。
\end{originaltheorem}
\paragraph*{注}
（a）若对偶性对不可约对成立，则对所有对成立。

（b）借助不变量理论的方法，可以证明此猜想“几乎”成立，见第11节。

（c）对 $G$ 或 $G'$ 各向异性（anisotropic）的各种对 $(G,G')$，许多作者已经研究过对偶猜想（例子见 [G]）。可以用（b）中的方法证明这些情形。对 $F=\R$，经典不变量理论的无穷小方法（infinitesimal method）可直接用于证明 $G$ 或 $G'$ 各向异性时的对偶性，见 [H 2]。

（d）若 $G$ 或 $G'$ 是 $\GL_1(D)$ 或 $O_2$，或者 $G=G'=T$ 为环面（torus），对偶性也成立。这些情形的理论大体相当于与广义理想类特征（Grössencharakter）相联系的 $L$ 函数的局部理论。

（e）对 $(\GL_n(F),\GL_n(F))$，对偶猜想本质上等价于 Weil [Wi 3] 关于 $\GL_n$ 在 $M_n$ 上左右作用的猜想。Weil 提出这一猜想，是为了给局部代数的 zeta 函数理论提供分布论的解释。
\origsection{7}{球面情形：Hecke 代数之间的对应}
$F$ 为剩余特征（residual characteristic）为奇数的非阿基米德局部域。

完整的局部对偶猜想可能很难证明，而自守形式理论中的某些应用只需较弱的结果。下面介绍一个与整体理论密切相关的局部定理。

令 $\mathcal O$ 为 $F$ 的整数环（integers）。设 $U$ 是 $F$ 上向量空间，配备形式 $(\ ,\ )$，可以是斜对称、对称或 Hermitian 形式（相对于 $F$ 的某个二阶自同构）。若 $S\subseteq U$ 是一个集合，令
\[
S^*=\{u\in U\mid(u,s)\in\mathcal O\text{，对所有 }s\in S\}.
\]
$U$ 中的格（lattice）指紧开 $\mathcal O$ 模。若 $L$ 为格，则 $L^*$ 也是格。若 $L=L^*$，称 $L$ 自对偶（self-dual）。若存在关于 $(\ ,\ )$ 的自对偶格，称该形式非分歧（unramified）。若 $G$ 为非分歧形式的等距群，则标准极大紧子群（standard maximal compact subgroup）指该形式的一个自对偶格的稳定子群（isotropy group）。在 $\GL_n(D)$ 中（$D$ 为未必在 $F$ 上中心的除代数），标准极大紧子群是任意在 $D$ 的整数环下不变的格的稳定子群。

$F$ 上的非分歧经典群（unramified classical group）指：
\begin{enumerate}
\item[(i)] $\GL_n(D)$，其中 $D$ 在 $F$ 上非分歧；或
\item[(ii)] $F'$ 上非分歧形式的等距群，其中 $F'/F$ 为非分歧扩张。
\end{enumerate}
回到辛向量空间 $W$。若不可约对 $(G,G')\subseteq\operatorname{Sp}$ 的两群 $G,G'$ 都非分歧，就称该对非分歧。一般的对非分歧，是指其每个不可约直和项都非分歧。

群 $\operatorname{Sp}$ 总是非分歧，而第3节定义的群 $J$（在那里记作 $J_v$）是 $\operatorname{Sp}$ 的标准极大紧子群。如第3节，可以无歧义地写 $\widetilde J=J\times\{1,-1\}$。若 $(G,G')$ 为 $\operatorname{Sp}$ 中的非分歧约化对偶对，$(K,K')$ 为 $(G,G')$ 的标准极大紧子群，则在共轭意义下可假定 $K,K'$ 包含于 $J$。于是可无歧义地写 $\widetilde K=K\times\{1,-1\}$，$K'$ 亦然。因此可以将 $K,K'$ 视为 $\widetilde{\operatorname{Sp}}$ 中明确给定的子群。
对具有标准极大紧子群 $K,K'$ 的非分歧对 $(G,G')$，令
\[
\mathcal R(\widetilde G,K)=\{\sigma\in\mathcal R(\widetilde G)\mid\sigma\text{ 有 }K\text{ 不变向量}\},
\]
对 $G'$ 以及 $G\cdot G'$ 也作类似定义。假定定义 $\omega_\chi$ 的特征 $\chi$ 非分歧，其意义为 $\mathcal O\subseteq\ker\chi\not\supseteq\varpi^{-1}\mathcal O$，其中 $\varpi$ 是 $F$ 的一个素元（prime）。

令 $C_c^\infty(\widetilde G//K)$ 为 $\widetilde G$ 上 $K$ 双不变函数（bi-invariant function）的（Hecke）代数；类似定义 $C_c^\infty(\widetilde G'//K')$。令 $I(K,K')$ 为同时被 $\omega_\chi^\infty(K)$ 和 $\omega_\chi^\infty(K')$ 固定的向量空间。显然 $\omega_\chi^\infty(C_c^\infty(\widetilde G//K))$ 保持 $I(K,K')$ 不变，因此可考虑其限制。
\begin{originaltheorem}{定理 7.1（Theorem 7.1）}
（a）$\mathcal R(\widetilde G\cdot\widetilde G',K\cdot K')$ 是 $\mathcal R(\widetilde G,K)$ 与 $\mathcal R(\widetilde G',K')$ 之间一个双射的图像。

（b）若 $\sigma\otimes\sigma'\in\mathcal R(\widetilde G\cdot\widetilde G')$，且 $\sigma\in\mathcal R(\widetilde G,K)$，则 $\sigma'\in\mathcal R(\widetilde G',K')$。

（c）限制
\[
\omega_\chi^\infty(C_c^\infty(\widetilde G//K))|_{I(K,K')},\quad
\omega_\chi^\infty(C_c^\infty(\widetilde G'//K'))|_{I(K,K')}
\]
为同一个算子代数。
\end{originaltheorem}
\origsection{8}{整体对偶性（Global duality）}
$F$ 现在又是整体域，相关对象仍按第3节定义。若 $G\subseteq\operatorname{Sp}$ 是 $F$ 上代数群，则可考虑 $\widetilde G_\A\subseteq\widetilde{\operatorname{Sp}}_\A$。仿照局部情形，定义
\[
\mathcal R(\widetilde G_\A)=\left\{\begin{array}{l|l}
\sigma&\sigma\text{ 是 }\widetilde G_\A\text{ 的不可约容许表示，且存在非平凡的}\\
&\widetilde G_\A\text{ 交织算子 }\alpha:\omega_\chi^\infty\longrightarrow\sigma
\end{array}\right\}.
\]
$\mathcal R(\widetilde G_\A)$ 中每个 $\sigma$，都在与 $\omega_\chi^\infty$ 相同的意义下是局部表示的张量积。容易看出，$\mathcal R(\widetilde G_\A)$ 本身是 $\mathcal R(\widetilde G_v)$ 关于紧子集 $\mathcal R(\widetilde G_v,K_v)$ 的限制直积，其中 $K_v=G\cap J_v$。
\begin{originaltheorem}{猜想：整体对偶性，第一部分（Conjecture: global duality, part I）}
若 $(G,G')\subseteq\operatorname{Sp}$ 为约化对偶对，则 $\mathcal R(\widetilde G_\A\cdot\widetilde G'_\A)$ 为 $\mathcal R(\widetilde G_\A)$ 与 $\mathcal R(\widetilde G'_\A)$ 之间一个双射的图像。
\end{originaltheorem}
这一猜想可直接由局部对偶猜想和第7节的定理推出。

设 $G\subseteq\operatorname{Sp}$ 如上。根据第3节的定理，可以把 $G$ 看作 $\widetilde G_\A$ 的子群。以 $\mathcal A(\widetilde G_\A)$ 表示 $G\backslash\widetilde G_\A$ 上缓慢增长的光滑函数。设 $\sigma$ 为 $\widetilde G_\A$ 的不可约容许表示。若存在 $\widetilde G_\A$ 嵌入 $\beta:\sigma\to\mathcal A(\widetilde G_\A)$，就称 $\sigma$ 自守（automorphic）。（这个自守表示概念比通常的概念 [P] 限制稍强。）对上述猜想的一个自然补充是：
\begin{originaltheorem}{猜想：整体对偶性，第二部分（Conjecture: global duality, part II）}
在上述 $\mathcal R(\widetilde G_\A)$ 与 $\mathcal R(\widetilde G'_\A)$ 的双射中，$\widetilde G_\A$ 的自守表示对应于 $\widetilde G'_\A$ 的自守表示，反之亦然。
\end{originaltheorem}
这个猜想同样有不变量理论的依据，这也把它与局部对偶猜想及第7节的定理联系起来。
\origsection{9}{超越的经典不变量理论：非阿基米德情形}
本节设 $F$ 为剩余特征为奇数的非阿基米德局部域，整数环为 $\mathcal O$。

下面说明第6--8节的猜想和结果与不变量理论之间的联系。为节省篇幅，只考虑非阿基米德域。这一情形与 [Wy] 中以代数方式表述的经典不变量理论仍有相当距离，因此两者之间的紧密联系不易直接看出。关于用约化对偶对重述经典不变量理论，包括在 $\R$ 上 $G$ 或 $G'$ 紧时局部对偶性的证明，见 [H 2]。

研究不变量理论，首先要构造一批自然的不变量。为此，回到辛向量空间 $W$ 及 Heisenberg 群 $H=H(W)$。令 $\chi$ 为 $H$ 的中心的非平凡特征。对任意子群 $A\subseteq H$，定义 $A^*$ 为模 $\ker\chi$ 意义下 $A$ 在 $H$ 中的中心化子。若 $A=A^*$，称 $A$ 极化（polarize）$\chi$。若 $A$ 极化 $\chi$，则 $F$（在此视为 $H$ 的中心）包含于 $A$，且 $A=(A\cap W)\oplus F$。集合 $A\cap W$ 为 $W$ 的一个子群，因此是 $\Z_p$ 模。若 $A\cap W$ 是 $\mathcal O$ 模，称 $A$ 为 $F$ 极化（$F$-polarization），或 $F$ 上的极化。以 $\Omega$ 表示所有极化的集合，$\Omega_F$ 为 $F$ 极化组成的子集。$\operatorname{Sp}$ 在 $H$ 上的作用置换这些极化。若 $\dim W=n$，则 $\Omega_F$ 由 $n+1$ 个 $\operatorname{Sp}$ 轨道组成。\footnote{\textbf{译者注（本次校订）：}这里的 $n$ 应为辛空间维数的一半，即 $\dim_F W=2n$。轨道按极化中向量子空间部分的维数 $0,\ldots,n$ 分类，共有 $n+1$ 个；下文 $\Omega_F^n$ 对应极大迷向子空间，也采用这一半维数。}

若 $A\subseteq H$ 极化 $\chi$，则存在 $\chi$ 到 $A$ 的唯一延拓 $\chi_A$，使 $A\cap W\subseteq\ker\chi_A$。
\begin{originaltheorem}{定理 9.1（Theorem 9.1）}
设表示 $\rho_\chi^\infty$ 实现在向量空间 $Y$ 上。对 $\chi$ 的每个极化 $A$，$Y$ 上都存在一个线性形式 $\lambda_A$，使
\[
\lambda_A(\rho_\chi^\infty(a)y)=\chi_A(a)\lambda_A(y),\qquad a\in A,\ y\in Y.\tag{9.1}
\]
方程 (9.1) 在相差常数倍的意义下确定泛函 $\lambda_A$。所得映射 $\mathcal Q:\Omega\to\mathbb P(Y^*)$ 将 $\Omega$ 嵌入为 $\mathbb P(Y^*)$ 的一个紧子集。对于 $\operatorname{Sp}$ 在 $\Omega$ 和 $\mathbb P(Y^*)$ 上的自然作用，映射 $\mathcal Q$ 等变（equivariant）。
\end{originaltheorem}
由于 $\mathbb P(Y^*)$ 有自然拓扑，可以通过嵌入 $\mathcal Q$ 赋予 $\Omega$ 拓扑。在此拓扑下，$\Omega_F$ 是 $\Omega$ 的闭子集，而 $\Omega_F$ 中每个 $\operatorname{Sp}$ 轨道带有其自然拓扑。可用 $\{\Omega_F^i\}_{i=0}^n$ 标记 $\Omega_F$ 中的 $\operatorname{Sp}$ 轨道，使 $\Omega_F^i$ 在 $\bigcup_{j=i}^n\Omega_F^j$ 中开且稠密。因此 $\Omega_F^n$ 是唯一闭轨道，由使 $A\cap W$ 为 $W$ 的极大迷向子空间的极化 $A$ 组成。$\Omega_F^0$ 则是唯一开轨道，由使 $A\cap W$ 为 $W$ 中的格（开紧 $\mathcal O$ 模）的 $A$ 组成。

下面利用 $\Omega_F$ 构造不变量。若 $G\subseteq\operatorname{Sp}$，$A\in\Omega_F$，且 $A$ 在 $G$ 下不变，则 $\lambda_A$ 必在相差常数倍意义下被 $\omega_\chi^{\infty *}(\widetilde G)$ 保持。因此 $\widetilde G$ 上存在拟特征（quasi-character）$\psi_A$，使
\[
\omega_\chi^{\infty *}(\widetilde g)\lambda_A=\psi_A(\widetilde g)\lambda_A,\qquad\widetilde g\in\widetilde G.\tag{9.2}
\]
因此，$\Omega_F$ 中的 $G$ 不动点给出 $Y$ 上的拟不变（quasi-invariant）线性泛函；在许多情形下，这些泛函实际上是不变的。这一构造与第4节利用 $\Theta$ 分布构造自守形式（$\theta$ 级数）的方法相对应。第12节将进一步发展这一类比。这在本质上也是经典不变量理论的第一基本定理（First Fundamental Theorem）的基础，其间的联系需要进一步解释。

设 $(G,G')\subseteq\operatorname{Sp}$ 为约化对偶对。若 $A\in\Omega_F^n$ 在 $G$ 下不变，令 $\psi_A$ 为 (9.2) 中与 $A$ 相联系的 $\widetilde G$ 上的拟特征。对任意 $g'\in G'$，$g'(A)$ 仍在 $G$ 下不变，且 $\psi_A=\psi_{g'(A)}$。因此，若以 $\Omega_F^n(G,\psi)$ 表示 $\Omega_F^n$ 中在 $G$ 下不变且 $\psi=\psi_A$ 的点 $A$ 的集合，则 $\Omega_F^n(G,\psi)$ 是 $G'$ 轨道之并。

由于 $\mathcal Q(\Omega_F^n(G,\psi))$ 为 $\mathbb P(Y^*)$ 的一个紧“子簇”，限制超平面截面丛（hyperplane section bundle），便定义线丛（line bundle）
\[
\begin{tikzcd}L(G,\psi)=L\arrow[d]\\\mathcal Q(\Omega_F^n(G,\psi)).\end{tikzcd}
\]
以 $\operatorname{Sec}^\infty(L)$ 表示 $L$ 的光滑截面空间。对 $\mathcal Q(\Omega_F^n(G,\psi))$ 上各点所代表的泛函求值，便得到线性映射 $e:Y\to\operatorname{Sec}^\infty(L(G,\psi))$。求值映射 $e$ 显然 $\widetilde G'$ 等变。其对偶为 $e^*:\operatorname{Sec}^{\infty *}(L)\to Y^*$。显然 $e^*$ 的像由 $\omega_\chi^{\infty *}(\widetilde G)$ 的特征泛函（eigenfunctional）组成，对应特征为 $\psi$。
\begin{originaltheorem}{定理 9.2（Theorem 9.2）}
若 $\Omega_F^n(G,\psi)$ 非空，则它由单个 $G'$ 轨道组成。此时 $L(G,\psi)$ 为 $\mathcal Q(\Omega_F^n(G,\psi))$ 上的 $\widetilde G'$ 齐性线丛（homogeneous line bundle）。最后，若 $\mu$ 为 $Y$ 上任意线性泛函，满足对 $\widetilde g\in\widetilde G$ 有 $\omega_\chi^{\infty *}(\widetilde g)\mu=\psi(\widetilde g)\mu$，则对某个 $\operatorname{Sec}^{\infty *}(L)$ 中的分布（distribution）$\nu$，有 $\mu=e^*(\nu)$。
\end{originaltheorem}
\origsection{10}{球面情形（The spherical case）}
$F$ 及其他对象同第9节。构成 $\Omega_F^0$ 的 $H$ 的子群模中心后为紧群；因此定理9.1、9.2 对 $\Omega_F^0$ 有更直接的类似结果。具体如下。
\begin{originaltheorem}{定理 10.1（Theorem 10.1）}
对每个 $A\in\Omega_F^0$，$Y$ 中存在在相差常数倍意义下唯一的向量 $y_A$，使 $\rho_\chi^\infty(a)y_A=\chi_A(a)y_A$（$a\in A$）。映射 $A\mapsto y_A$ 对于 $\operatorname{Sp}$ 在 $\Omega_F^0$ 与 $\mathbb P(Y)$ 上的相应作用等变。
\end{originaltheorem}
设 $(G,G')$ 是 $\operatorname{Sp}$ 中的非分歧约化对偶对，$(K,K')$ 是一对标准极大紧子群，二者都包含于 $\operatorname{Sp}$ 的标准极大紧子群 $J$。如第7节所述，可以将 $J$ 从而将 $K,K'$ 看作嵌入 $\widetilde{\operatorname{Sp}}$ 的子群。令 $L\subseteq W$ 为 $J$ 固定的自对偶格。假定 $\chi$ 在第7节的意义下非分歧，则 $A=L\oplus F\subseteq H$ 极化 $\chi$。由于 $J$ 正规化 $A$，向量 $y_A$ 在相差常数倍的意义下被 $\omega_\chi^\infty(J)$ 保持。因为 $J$ 是完美群（perfect group），可见 $y_A$ 实际上在 $\omega_\chi^\infty$ 下被 $J$ 固定，从而也被 $K,K'$ 固定。（例外是 $F$ 的剩余域有三个元素且 $\dim W=2$，这一情形可单独处理，但在此排除。）
\begin{originaltheorem}{定理 10.2（Theorem 10.2）}
对 $f\in C_c^\infty(\widetilde G'/K')$，定义 $r(f)=\omega_\chi^\infty(f)(y_A)$。则 $r$ 是从 $C_c^\infty(\widetilde G'/K')$ 到 $Y$ 中 $K$ 不变向量空间 $I(K)$ 的满射。
\end{originaltheorem}
由于 $\widetilde G'$ 的每个具有 $K'$ 不变向量的不可约容许表示，都在相差非零标量的意义下唯一地成为 $C_c^\infty(\widetilde G'/K')$ 的商，又由于 $C_c^\infty(\widetilde G'/K')^{K'}=C_c^\infty(\widetilde G'//K')$ 是有单位元的交换代数（因而作为自身的模是循环模），定理7.1 可由定理10.2 推出。\footnote{\textbf{译者注（本次校订）：}本段原来引用“定理9.4”，但本文没有这一编号；所用的是紧接上文的定理10.2。}
\origsection{11}{变量加倍与局部对偶性（Doubling variables and local duality）}
$F$ 及其他对象继续同第9、10节。

下面说明如何利用定理9.2 研究局部对偶猜想。对 $\widetilde G$ 的不可约容许表示 $\sigma$，令 $Y_\sigma$ 为 $Y$ 对所有 $\widetilde G$ 交织映射 $Y\to\sigma$ 的核的交所取的商。显然 $\omega_\chi^\infty(\widetilde G)$ 在 $Y_\sigma$ 上诱导的表示是 $\sigma$ 的若干副本之和。也不难看出，从 $Y$ 到 $Y_\sigma$ 的商映射的核在 $\widetilde G'$ 下不变，因此 $Y_\sigma$ 实际是 $\widetilde G\times\widetilde G'$ 模。作为这种模，它必形如 $\sigma\otimes\tau$，其中 $\tau$ 是 $\widetilde G'$ 的某个光滑模。$\mathcal R(\widetilde G\cdot\widetilde G')$ 为从 $\mathcal R(\widetilde G)$ 到 $\mathcal R(\widetilde G')$ 的单射这一条件（术语同第6节），等价于说：只要 $Y_\sigma\ne\{0\}$，表示 $\tau$ 就有唯一的非零不可约容许 $\widetilde G'$ 模 $\sigma'$ 作为商。

考虑 $\sigma'$ 的唯一性问题，它比存在性更为关键。假定 $\sigma$ 的（光滑）逆反表示（contragredient）$\sigma^*$ 为 $\omega_{\bar\chi}^\infty$ 的商，其中 $\bar\chi=\chi^{-1}$ 是 $F$ 上与 $\chi$ 互逆的特征。如果 $\sigma$ 是酉的，甚至只是拟酉的（quasi-unitary），这一点便自动成立。令 $Y^\vee$ 为 $\omega_{\bar\chi}^\infty$ 的空间，$Y_{\sigma^*}^\vee$ 为 $Y^\vee$ 的最大商，使 $\widetilde G$ 在其上以 $\sigma^*$ 的重数倍作用。于是 $Y_{\sigma^*}^\vee$ 是形如 $\sigma^*\otimes\tau^\vee$ 的 $\widetilde G\times\widetilde G'$ 模，其中 $\tau^\vee$ 是另一个 $\widetilde G'$ 模。因此 $Y_\sigma\otimes Y_{\sigma^*}^\vee$ 是形如 $(\sigma\otimes\sigma^*)\otimes(\tau\otimes\tau^\vee)$ 的 $(\widetilde G\times\widetilde G)\times(\widetilde G'\times\widetilde G')$ 模。

令 $\widetilde G_\Delta$ 为 $\widetilde G\times\widetilde G$ 的对角子群（diagonal subgroup）。那么 $\sigma\otimes\sigma^*$ 限制到 $\widetilde G_\Delta$，恰以唯一方式具有平凡表示作为商。因此，若 $(Y\otimes Y^\vee)_1$ 是 $Y\otimes Y^\vee$ 的使 $\widetilde G_\Delta$ 平凡作用的最大商，则 $\tau\otimes\tau^\vee$ 为 $(Y\otimes Y^\vee)_1$ 的一个 $\widetilde G'\times\widetilde G'$ 商。因此，若能分析 $(Y\otimes Y^\vee)_1$ 作为 $\widetilde G'\times\widetilde G'$ 模的结构，特别是证明它仅有形如 $\sigma'\otimes\sigma'^*$ 的不可约商（$\sigma'$ 为 $\widetilde G'$ 的不可约容许表示），就得到局部对偶猜想的一个接近目标的版本。事实上，梳理以上化约，可以归纳出三个左右的技术性猜想，它们合起来与局部对偶猜想等价。

上述化约的要点在于，可以利用定理9.2 描述 $(Y\otimes Y^\vee)_1$。具体来说，令 $W^-$ 为底层空间是 $W$、辛形式为 $-\langle\ ,\ \rangle$（即 $W$ 上形式之负）的辛向量空间。为区分这两个辛结构，将原来的 $W$ 记作 $W^+$。令辛向量空间 $2W=W^+\oplus W^-$，称为 $W$ 的加倍（double）。按第1节的办法构造 $H(2W)$。可由公式
\[
\begin{aligned}
i_1(w,t)&=((w,0),t),\\
i_2(w,t)&=((0,-w),-t),
\end{aligned}\qquad w\in W,\ t\in F,\tag{11.1}
\]
定义 $H(W)$ 到 $H(2W)$ 的单射 $i_1,i_2$。因此 $i_1\times i_2:H(W)\times H(W)\to H(2W)$ 为满同态，其核为中心的对角子群。令 $\chi$ 为 $F$ 的特征，$2\rho_\chi$ 为 $H(2W)$ 的相应表示。容易看出
\[
2\rho_\chi\circ(i_1\times i_2)\cong\rho_\chi\otimes\rho_{\bar\chi}\tag{11.2}
\]
（外张量积，outer tensor product）。

令 $\operatorname{Sp}(2W)$ 为 $2W$ 的辛群，$\operatorname{Sp}(W)$ 为 $W$ 的辛群。有两个从 $\operatorname{Sp}(W)$ 到 $\operatorname{Sp}(2W)$ 的嵌入，与 (11.1) 的映射相容：
\[
\begin{aligned}
i_1(g)(w_1,w_2)&=(g(w_1),w_2),\\
i_2(g)(w_1,w_2)&=(w_1,g(w_2)),
\end{aligned}\quad g\in\operatorname{Sp}(W),\ (w_1,w_2)\in2W.\tag{11.3}
\]
这些映射唯一提升为相应 $\widetilde{\operatorname{Sp}}$ 群之间的映射。令 $2\omega_\chi$ 为由 (2.1) 定义的 $\widetilde{\operatorname{Sp}}(2W)$ 的表示。与 (11.2) 一样，容易看出
\[
2\omega_\chi\circ(i_1\times i_2)=\omega_\chi\otimes\omega_{\bar\chi}\tag{11.4}
\]
（外张量积）。

设 $(G,G')$ 为 $\operatorname{Sp}(W)$ 中的约化对偶对。可以证明存在群 $2G'$，使 $(i_1\times i_2(G_\Delta),2G')$ 构成 $\operatorname{Sp}(2W)$ 中的约化对偶对。群 $2G'$ 与 $G'$ 的关系，如同 $\operatorname{Sp}(2W)$ 与 $\operatorname{Sp}(W)$ 的关系；特别地，$i_1(G')\times i_2(G')\subseteq2G'$。

从这一加倍构造可见，以上定义的空间 $(Y\otimes Y^\vee)_1$ 作为 $2\widetilde G'$ 模，可由定理9.2 描述。通过限制，便能考察其作为 $\widetilde G'\times\widetilde G'$ 模的结构。由此可得：若 $G$ 或 $G'$ 紧，对偶猜想确实成立；一般情形则“几乎”成立。一般情形还需要补上剩余的论证。有限域情形的更多细节见 [H 1]。
\origsection{12}{整体不变量（Global invariants）}
本节中，$F$ 为整体域，相关结构同第3节。给定 $F$ 模 $M\subseteq W_\A$，令
\[
M^*=\{x\in W_\A\mid\langle x,m\rangle\in F\text{，对所有 }m\in M\}.
\]
若 $M=M^*$，称 $M$ 为 $W$ 的整体极化（global polarization）。每个整体极化 $M$ 都对应一个闭子群 $M\oplus F\subseteq H_\A(W)$。下面与定理9.1、10.1 对应的整体结果，推广了定理4.1。以 $\Omega_\A$ 表示整体极化空间。
\begin{originaltheorem}{定理 12.1（Theorem 12.1）}
设 $H_\A$ 的表示 $\rho_\chi$ 实现在空间 $Y$ 上。\footnote{\textbf{译者注（本次校订）：}此处用于定义分布的 $Y$ 应沿用第4节的光滑向量空间，即 $\rho_\chi^\infty$ 的空间。$\Theta$ 这类求和分布一般不是原 Hilbert 空间上的有界泛函。}对每个 $M\in\Omega_\A$，$Y$ 上存在在相差常数倍意义下唯一的线性泛函 $\theta_M$，使
\[
\rho_\chi^{\infty *}(x)\theta_M=\theta_M,\qquad x\in M\oplus F\subseteq H_\A.
\]
所得映射 $\mathcal Q:\Omega_\A\to\mathbb P Y^*$ 是 $\Omega_\A$ 的 $\operatorname{Sp}_\A$ 等变嵌入，且像紧。
\end{originaltheorem}
显然 $W_F\subseteq W_\A$ 是整体极化，且 $\mathcal Q(W_F)=\Theta$ 即第4节的 $\theta$ 分布。$W_F$ 在 $\operatorname{Sp}_\A$ 中的稳定子群为 $\operatorname{Sp}_F$，因此 $\Theta$ 的 $\operatorname{Sp}_\A$ 轨道同构于 $\operatorname{Sp}_\A/\operatorname{Sp}_F$，它是 Siegel 上半空间（Siegel upper half-plane）的整体版本。此外，$\Theta$ 的 $\operatorname{Sp}_\A$ 轨道是 $\mathcal Q(\Omega_\A)$ 中唯一的开轨道且稠密，所以 $\Omega_\A$ 是 $\operatorname{Sp}_\A/\operatorname{Sp}_F$ 的紧化（compactification）。嵌入 $\mathcal Q$ 是 Igusa [I] 研究的 Siegel 模簇（Siegel modular variety）的射影嵌入的整体类似物，而且后者可由它导出。F. Herman 正在研究这一问题。从这一观点研究 Shimura [Sh] 的典范模型（canonical model），也可能得到进一步结果。

如第9节，$\mathbb P Y^*$ 上的超平面截面丛在 $\mathcal Q(\Omega_\A)$ 上诱导线丛 $L$，并有从 $Y$ 到 $L$ 的截面的求值映射。这些截面正是第4节及经典理论中的 $\theta$ 级数。（由于 $\Theta$ 实际在 $\operatorname{Sp}_F$ 下不变，丛 $L$ 在 $\Theta$ 的轨道上是 $\widetilde{\operatorname{Sp}}_\A$ 等变平凡的，而取 $\Theta$ 的 $\widetilde{\operatorname{Sp}}_\A$ 轨道，便产生 $L$ 的一个 $\widetilde{\operatorname{Sp}}_\A$ 不变截面。$L$ 的截面限制到这一横截面，就得到 $\theta$ 级数。）$\theta$ 级数可延拓为整个 $\mathcal Q(\Omega_\A)$ 上的截面，保证其在 $\operatorname{Sp}_\A/\operatorname{Sp}_F$ 的“无穷远处”有良好行为。

设 $(G,G')$ 为 $\operatorname{Sp}_F$ 中的约化对偶对。$\theta$ 分布在 $\operatorname{Sp}_F$ 下不变，因而在 $G$ 下不变。所以 $\Theta$ 在 $\mathcal Q(\Omega_\A)$ 中的整个 $\widetilde G'_\A$ 轨道都由 $G$ 不变泛函组成。类比定理9.1（它本身又类比于经典不变量理论），自然会进一步问，这种方式能否给出全部 $G$ 不变泛函。相应的弱形式可表述为以下猜想。
\begin{originaltheorem}{猜想：弱整体不变量（Conjecture: weak global invariants）}
令 $Y_1\subseteq Y$ 为 $Y$ 上所有 $G_F$ 不变泛函的核之交，$Y_2$ 为所有 $g'\in\widetilde G'_\A$ 对应的 $\omega_\chi^{\infty *}(g')\Theta$ 的核之交。则 $Y_2=Y_1$。
\end{originaltheorem}
若 $G_F'\backslash\widetilde G'_\A$ 紧，则此猜想蕴含：$\mathcal R(\widetilde G_\A)$ 中每个自守表示，都可以通过取 $\widetilde G_\A\cdot\widetilde G'_\A$ 上的 $\theta$ 级数，再与 $G'_F\backslash\widetilde G'_\A$ 上的某个函数相乘积分来得到。特别是，弱整体不变量猜想与第一个整体对偶猜想，将蕴含：若 $\sigma\in\mathcal R(\widetilde G_\A)$ 自守，则 $\sigma'\in\mathcal R(\widetilde G'_\A)$ 也自守，因而证明第二个整体对偶猜想的一个方向。

仍在 $G'_F$ 各向异性的情形下，整体不变量猜想的另一观点是，它给出一个局部整体原理（local-global principle）：若 $\sigma$ 为自守表示，且其每个局部分量都是局部振子表示的商，则 $\sigma$ 可由 $\theta$ 级数得到。由此可见，$\theta$ 级数之间对应关系的主要问题，都与不变量理论紧密相连。

\begin{thebibliography}{Wi 3}
\bibitem[C1]{c17-c1} P. Cartier, \emph{Quantum mechanical commutation relations and theta functions}, Proc. Sympos. Pure Math., vol. 9, Amer. Math. Soc., Providence, R.I., 1966, pp. 361--383.
\bibitem[C2]{c17-c2} P. Cartier, \emph{Representations of $p$-adic groups}, 本论文集，第1册，111--155页。
\bibitem[F]{c17-f} D. Flath, \emph{Decomposition of representations into tensor products}, 本论文集，第1册，179--183页。
\bibitem[G]{c17-g} S. Gelbart, \emph{Examples of dual reductive pairs}, 本论文集，第1册，287--296页。
\bibitem[H1]{c17-h1} R. Howe, \emph{Invariant theory and duality for classical groups over finite fields}, 预印本。
\bibitem[H2]{c17-h2} R. Howe, \emph{Remarks on classical invariant theory}, 预印本。
\bibitem[I]{c17-i} J. Igusa, \emph{Theta functions}, Die Grundlehren der Math. Wissenschaften, Band 194, Springer-Verlag, Berlin, 1972.
\bibitem[M]{c17-m} C. Moore, \emph{Group extensions of $p$-adic linear groups}, Inst. Hautes Etudes Sci. Publ. Math. 35 (1968).
\bibitem[P]{c17-p} I. I. Piatetski-Shapiro, \emph{Classical and adelic automorphic forms. An introduction}, 本论文集，第1册，185--188页。
\bibitem[Sa]{c17-sa} I. Satake, \emph{Symplectic representations of algebraic groups satisfying a certain analyticity condition}, Acta Math. 117 (1961), 215--279.
\bibitem[Sh]{c17-sh} G. Shimura, \emph{Moduli of abelian varieties and number theory}, Proc. Sympos. Pure Math., vol. 9, Amer. Math. Soc., Providence, R.I., 1966, pp. 312--332.
\bibitem[Si 1]{c17-si1} C. L. Siegel, \emph{Indefinite quadratische Formen und Funktionen Theorie. I}, Math. Ann. 124 (1951), 17--24; II, 124 (1952), 364--387.
\bibitem[Si 2]{c17-si2} C. L. Siegel, \emph{Über die Zeta functionen indefiniter quadratischen Formen}, Math. Z. 43 (1938), 682--708.
\bibitem[Wa]{c17-wa} N. Wallach, \emph{Representations of reductive Lie groups}, 本论文集，第1册，71--86页。
\bibitem[Wi 1]{c17-wi1} A. Weil, \emph{Sur certains groupes d'opérateurs unitaires}, Acta Math. 111 (1964), 143--211.
\bibitem[Wi 2]{c17-wi2} A. Weil, \emph{Sur la formule de Siegel dans la théorie des groupes classiques}, Acta Math. 113 (1965), 1--87.
\bibitem[Wi 3]{c17-wi3} A. Weil, \emph{Fonction zeta et distributions}, Sem. Bourbaki, exposé 312, June 1966.
\bibitem[Wy]{c17-wy} H. Weyl, \emph{The classical groups}, Princeton Univ. Press, Princeton, N.J., 1946.
\end{thebibliography}
\noindent 耶鲁大学（Yale University）。
\par\noindent\footnotesize 原文 AMS（MOS）学科分类（1970）：主分类 22E55。版权 © 1979 American Mathematical Society。\normalsize
